\section{Limitations and Open Questions}
\label{sec:limitations}

The evidence is simulation-only and limited to the related pick-and-place tasks
of LIBERO Object, one frozen ACT checkpoint, and one partition: six
arm/end-effector dimensions versus a one-dimensional gripper.  It does not
establish cross-suite, cross-policy, cross-embodiment, dexterous-hand, or
real-robot generalization.  RoboTwin, RoboDojo, force sensing, and physical-robot
evaluation are pending directions rather than reported experiments.

All evaluated horizons are static.  Per-task best configurations are
retrospective oracle-style analyses and cannot be deployed without a selection
rule.  The strong common $(4,16)$ baseline further means that current results do
not demonstrate a need for within-episode adaptation.  A next-stage study must
separate task-level static selection from genuine phase-dependent gain, define
an online observable signal, and compare at controlled query budgets.

Finally, group-specific execution deliberately composes actions from different
source chunks.  The retained components were jointly predicted with components
that may no longer be executed, so temporal inconsistency can arise across
groups.  We have not established why particular horizon assignments succeed.
Exploratory trace differences in action variation and sign changes are
descriptive, not causal.  Future work should test synchronization rules and
cross-chunk consistency controls before learning a group-specific scheduler.
